\chapter*{Abstract}
\addcontentsline{toc}{chapter}{Abstract}

The rise of distance learning raises growing concerns about the integrity of assessments, and this is precisely the problem this project aims to address. Certif-fun is a platform that brings together in a single tool the management of certifications, automated exam proctoring, and pedagogical support. The architecture relies on distributed microservices combining Java Spring Boot, React, and FastAPI. For proctoring, computer vision models detect in real time suspicious gaze patterns and unauthorized objects, while identity verification through QR code scanning and facial recognition secures access to exams. Automatic quiz generation is handled by a local AI via Ollama. Tests carried out show that the system correctly identifies suspicious behaviors, prolonged absences, and continuous background noise, with a response time suitable for real exam conditions. The interface was rated smooth and non-intrusive by tested users. Certif-fun demonstrates that it is possible to combine strict security and user comfort within a single tool, paving the way for online certifications whose credibility can rival that of in-person exams.

\bigskip
\noindent \textbf{Keywords:} Online Certification, Automated Proctoring, Assessment Integrity, Artificial Intelligence, Microservices, Computer Vision.
